\subsection{The model}\label{\refsec:model}
We consider an economy populated by overlapping generations of workers and retirees. Within each generational cohort, we consider $J+1$ different household types indexed $j= 0,...,J$. We refer to type $j=0$ households as "informal"; they work in informal markets and save through an informal technology that is not part of the formal capital stock. Households of type $i>0$ are "formal" households; they supply $h_{t,i}\eta_{t,i}$ units of labour when young with $\eta_{t,i}$ denoting labour productivity. We let $\gamma_{t,j}$ denote sizes of working households with $\sum_i \gamma_{t,i} = 1$.

As young, households are endowed with a unit of time that is allocated between labour and leisure. Formal workers pay taxes on their labour income, consume, and save; informal workers consume and save through an informal technology. Households survive and retire with probability $p_{t,j}$; all formal households survive with the same probability $p_{t,i}= p_t$. Lost savings are insured within-type, for informal as well as formal households.

The government runs a pay-as-you-go (PAYG) pension system with a balanced budget. Each period, the government imposes a tax on labour supplied in the formal sector and transfers the collected amount directly to retirees. We consider social security systems that differ along two dimensions: whether workers need to qualify to receive full benefits or not, and the type of benefits provided - either tied to earnings or flat-rate.

Policy decisions are taken by a government that acts in the interest of voters, but lacks commitment.

\smalltitle{Technology.}
\begin{subequations}\label{\refeq:aggregateStates}
In the formal sector, a continuum of firms transforms aggregate capital and labour into output. We define aggregate labour and capital from
\begin{align}
    H_t =& n_th_t, & h_t&\equiv \sum_{i}\gamma_{t,i}\eta_{t,i}h_{t,i} \\
    K_t =& n_{t-1} s_{t-1}, & s_{t-1}&\equiv \sum_i\gamma_{t-1,i}s_{t-1,i}.
\end{align}
\end{subequations}
Both sums run over the formal types $i>0$ only: informal labour is supplied outside the formal sector, and informal savings $s_{t,0}$ are held in the informal savings vehicle and do \emph{not} add to the formal capital stock. The representative firm produces output using labour and capital and a conventional Cobb-Douglas technology
\begin{align*}
    Y_t = A_t K_t^{\alpha}H_t^{1-\alpha}.
\end{align*}
In equilibrium this implies the following factor prices:
\begin{align}\label{\refeq:factorPrices}
    R_t = \alpha \left(\dfrac{s_{t-1}}{\nu_t h_t}\right)^{\alpha-1}, \qquad w_t = (1-\alpha)A_t \left(\dfrac{s_{t-1}}{\nu_t h_t}\right)^{\alpha},
\end{align}


\bigskip
Informal households employ a linear production technology with return $w_t^0$ per efficiency unit of labour and $R_t^0$ per unit of informal savings. For simplicity, we model both as proportional to the corresponding formal factor price absent taxes writing
\begin{align}\label{\refeq:factorPrices0}
    w_t^0= \left(\frac{1-\alpha}{\Gamma_{h,t}^{\alpha}}\right)^{\frac{1}{1+\alpha\xi}} \left(\frac{s_{t-1}}{\nu_t}\right)^{\frac{\alpha}{1+\alpha\xi}}, \qquad R_t^0= \alpha \left(\frac{s_{t-1}}{\nu_th_t}\right)^{\alpha-1} \chi_t^R,
\end{align}
where $\chi_t^R$ scales the return relative to the formal capital market. We note that a similar scaling is not required for $w_t^0$ as the informal income already includes the type-specific productivity $\eta_{t,0}$. 


\smalltitle{Preferences and Household Choices.}
Workers and retirees in period $t$ value consumption, $c_{1,t}^j$ and $c_{2,t}^j$ respectively. Workers discount the future at factor $\beta_{t,j}\in[0,1)$ where the survival probability $p_{t,j}$ is included.

For the young households, the expected utilities are given by
\begin{align}\label{\refeq:util}
    &\max\mbox{ } U_{t,j}\left(c_{1,t}^j, h_{t,j}, c_{2,t+1}^j\right) = u_{t,j}\left(c_{1,t}^j, h_{t,j}\right) +\beta_{t,j}u_{t,j}\left(c_{2,t+1}^j,0\right)
\end{align}
where we use CRRA-GHH preferences
\begin{align*}
    u_{t,j}(c,h) = \left\lbrace\begin{array}{ll} \dfrac{1}{1-1/\rho}\left(\tilde{c}_{t}^j\right)^{1-1/\rho}, & \text{for }\rho \neq 1 \\
    \ln\left(\tilde{c}_{t}^j\right) , & \text{for }\rho = 1
    \end{array}\right. , \qquad \tilde{c}_t^j \equiv c-\dfrac{\xi}{1+\xi}X_{t,j}\left(h\right)^{\frac{1+\xi}{\xi}}
\end{align*}
where $\rho>0$ is the intertemporal elasticity of substitution, $\xi$ is the Frisch elasticity, and $X_{t,j}$ is a parameter that measures the disutility of labour.

\begin{subequations}\label{\refeq:formalBudget}
The formal household receives wages, pay social security taxes $\tau_t$, save $s_{t,i}$, and consume when young, and consume their savings including pension benefits $b_{t+1}^i$ when old:
\begin{align}
    c_{1,t}^i &= w_t\eta_{t,i} h_{t,i}(1-\tau_t)-s_{t,i} \\
    c_{2,t+1}^i &= s_{t,i}R_{t+1}/p_t+ b_{t+1}^i(h_{t,i}).
\end{align}
\end{subequations}

\begin{subequations}\label{\refeq:informalBudget}
Informal households earn informal wages when young, of which they consume part and save part in the informal vehicle; when old they consume their informal savings including pension benefits $b_{t+1}^0$:
\begin{align}
    c_{1,t}^0 &= w_t^0\eta_{t,0}h_{t,0} - s_{t,0} \\
    c_{2,t+1}^0 &= s_{t,0}R^0_{t+1}/p_{t,0}+b_{t+1}^0,
\end{align}
As for formal households, savings lost to mortality are insured within the type, which is the source of the $1/p_{t,0}$ term.\footnote{Note that we do not formally restrict $s_{t,0}$ to be non-negative. In all of our calibrations, however, $s_{t,0}/s_t$ has turned out to be safely inside the unit interval.}
\end{subequations}



\smalltitle{Social Security System.}
The government runs a pay-as-you-go pension system with a balanced budget. Every period, the government taxes labour income in the formal sector and transfers the sum directly to current retirees. Informal savings are not taxed.

Benefits are given by
\begin{subequations} \label{\refeq:governmentBudget}
\begin{align}
    b_t^i =& \left[\theta_th_{t-1,i} \eta_{t-1,i}+(1-\theta_t)h_{t-1}\right] \overline{b}_t \\
    b_t^0 =& \epsilon_t h_{t-1} \overline{b}_t
\end{align}
where $\epsilon_t$ measures universal pensions and $\theta_t$ measures contributive incentives.



We note that the shares $\sum_j\gamma_{t,j} = 1+\gamma_{t,0}$ to have a unit mass of formal households (convenient to have for definition of aggregate states). Thus, the average contribution of young households at time $t$ is
\begin{align*}
    \text{Avg. contribution from formal} = w_t\tau_t\sum_{i>0}\gamma_{t,i}\eta_{t,i}h_t^i = w_t\tau_th_t.
\end{align*}
Letting $n_t$ denote total number of young households, the number of formal households are $n_t/(1+\gamma_{t,0})$. Similarly, without mortality, the share of retired households of type $j$ would be $\gamma_{t-1,j}/(1+\gamma_{t-1,0})$. Taking mortality into consideration, the balanced budget requires that
\begin{align*}
    \dfrac{n_t}{1+\gamma_{t,0}} w_t\tau_th_t = \dfrac{n_{t-1}}{1+\gamma_{t-1,0}}\sum_j \gamma_{t-1,j}p_{t-1,j} b_t^j
\end{align*}
The left hand side measures contributions, the right hand side measures beneficiaries. Using the structure of the pension design, this identifies the level in pension rates $\overline{b}_t$ as:
\begin{align}\label{\refeq:governmentBudget:bbar}
    \overline{b}_t & \equiv \dfrac{\nu_t w_th_t\tau_t}{h_{t-1}\kappa_{t-1}}, \qquad\kappa_t \equiv \frac{1+\gamma_{t,0}}{1+\gamma_{t+1,0}}(\epsilon_{t+1}\gamma_{t,0}p_{t,0}+p_t)
\end{align}
\end{subequations}


\bigskip
\begin{subequations}\label{\refeq:formalOpt}
With CRRA-GHH preferences, households maximize \eqref{\refeq:util} subject to budgets \eqref{\refeq:formalBudget} and \eqref{\refeq:informalBudget}. The labour supply and savings of formal households are then given by (given factor prices and policies):
\begin{align}
    h_{t,i} &= \left[\frac{\eta_{t,i}}{X_{t,i}}\left(w_t(1-\tau_t)+\theta_{t+1}\overline{b}_{t+1}\right)\right]^{\xi} \\
    s_{t,i} &= \dfrac{B_{t+1}^i}{1+B_{t+1}^i}\left[w_t\left(1-\tau_t\right)h_{t,i}\eta_{t,i}-\dfrac{X_{t,i}\xi}{1+\xi}\left(h_{t,i}\right)^{\frac{1+\xi}{\xi}}\right]-\dfrac{b_{t+1}^i/(R_{t+1}/p_t)}{1+B_{t+1}^i}
\end{align}
\end{subequations}
where we have used the shorthand $B_{t+1}^i \equiv \left(\beta_{t,i}\right)^{\rho}(R_{t+1}/p_t)^{\rho-1}$.

For the informal households, similar conditions apply with a few adjustments for relevant factor prices and the incentive structure of universal pensions:
\begin{subequations}\label{\refeq:informalOpt}
\begin{align}
    h_{t,0} &= \left(\frac{\eta_{t,0}}{X_{t,0}}w_t^0\right)^{\xi},\\
    s_{t,0} &=  \dfrac{B_{t+1}^0}{1+B_{t+1}^0}\left[w_t^0h_{t,0}\eta_{t,0}-\dfrac{X_{t,0}\xi}{1+\xi}\left(h_{t,0}\right)^{\frac{1+\xi}{\xi}}\right]-\dfrac{b_{t+1}^0/(R_{t+1}^0/p_t^0)}{1+B_{t+1}^0},
\end{align}
\end{subequations}
where $B_{t+1}^0 \equiv \left(\beta_{t,0}\right)^{\rho} \left(R_{t+1}^0/p_{t,0}\right)^{\rho-1}$ is the generalized discount function for the informal households.


